\documentclass[10pt,a4paper]{amsart}
\usepackage[margin=19mm]{geometry}
\usepackage{amsmath,amssymb,mathtools}
\usepackage[scr=rsfs,cal=euler]{mathalfa}
\usepackage{xfrac}
\usepackage[unicode,hidelinks,bookmarks=false]{hyperref}
\newcommand{\ie}{i.e.\ }
\newcommand{\cf}{cf.\ }
\newcommand{\resp}{respectively\ }
\newcommand{\Subsection}{\subsection}
\providecommand{\nobreakdash}{\nobreak\hbox{-}}
\setlength{\emergencystretch}{2em}
\multlinegap0pt
\usepackage{bm}
\usepackage{bbm}
\usepackage{dsfont}
\usepackage{xspace}
\usepackage{cancel}
\usepackage{mathtools}
\usepackage{amsthm}
\makeatletter
\@ifundefined{theorem}{\newtheorem{theorem}{Theorem}}{}
\@ifundefined{lemma}{\newtheorem{lemma}[theorem]{Lemma}}{}
\@ifundefined{proposition}{\newtheorem{proposition}[theorem]{Proposition}}{}
\makeatother
\newcommand{\ch}{\mathrm{ch}}
\newcommand{\single}{\chi_{\nleftrightarrow}}
\title[Single conflict coloring, adaptable choosability and separat]{Single conflict coloring, adaptable choosability and separation choosability}
\author{Carl Johan Casselgren, Kalle Eriksson}
\date{Source: arXiv:2509.13913v1 (CC BY 4.0)}
\begin{document}
\maketitle

\noindent\emph{Source and licence.} Single conflict coloring, adaptable choosability and separation choosability. Version arXiv:2509.13913v1 (CC BY 4.0), distributed under the Creative Commons Attribution 4.0 International licence, as recorded in the arXiv licence metadata for this exact version and shown on the abstract page https://arxiv.org/abs/2509.13913v1. Full rights notice: licenses/arxiv-2509.13913-CC-BY-4.0.txt.\par\smallskip

\noindent\textit{Editorial scope.} Counted unit: the Theorem whose source label is \texttt{prop:noplanar} in the article (source0.tex lines 867--871, source label \texttt{prop:noplanar}), delivered together with the article's own complete original proof of it (source0.tex lines 872--961, one \texttt{proof} environment of 90 lines). No mathematics is written, repaired or re-derived: statement and proof are the article's own text line for line. Required context retained uncounted: prop:disjointcycles (lines 414--420); prop:maxdeg (lines 508--511); lem:wheel (lines 859--862). Every definition, display and label of those lines is retained unchanged. Omitted from this package and not redistributed: 1--413, 421--507, 512--858, 863--866, 962--1287. Nothing inside a retained excerpt is omitted or altered: every retained body is delivered exactly as the article's lines are written, so no omission receipt of any kind accompanies this package. Citations: the retained excerpts carry no citation command at all, so no bibliography entry of the article is transcribed and no reference is redistributed. Reference adaptation: none was needed and none was made; every reference inside the delivered unit resolves inside it, so no unresolved reference or citation is produced. Printed slips kept literal: no printed word, symbol, hypothesis or number of the counted statement or of its proof was altered. Layout and class: the article's own class is \texttt{article} with packages including inputenc, amssymb, geometry, amsmath, amsthm, graphicx, float, longtable, xcolor, cleveref, subcaption, enumitem, dsfont, caption; the wrapper is the lane's pinned \texttt{amsart} editorial wrapper at 10pt on A4, which restates the theorem-like environments the retained text needs and adds the article's own macro definitions that the retained text needs (\texttt{\textbackslash ch}, \texttt{\textbackslash single}), with extra wrapper packages (bm, bbm, dsfont, xspace, cancel, mathtools). The delivered numbering is the unit's own. Assets: no retained span contains a figure environment, a graphics-inclusion command, a PDF-page inclusion, a table or a TikZ picture; no image file is copied into this package, the package declares no asset dependency and no image file is redistributed with it. Licence: version arXiv:2509.13913v1 of this article is distributed under the Creative Commons Attribution 4.0 International licence, as recorded in the arXiv licence metadata for this exact version and shown on the abstract page https://arxiv.org/abs/2509.13913v1. A full rights notice is delivered at licenses/arxiv-2509.13913-CC-BY-4.0.txt. Characters: every retained excerpt is pure ASCII, so the delivered engine, XeLaTeX with its default Latin Modern text font, reports no missing character.
\begin{proposition}
\label{prop:disjointcycles}
If a connected graph has (at least) two
cycles of length at least $4$ with at most one common vertex,
 then it is 
not separation $2$-choosable.
\end{proposition}

\begin{proposition}
\label{prop:maxdeg}
	For any graph $G$, $\single(G) \leq \lceil \Delta(G)/2 \rceil +1$.
\end{proposition}

\begin{lemma}
    \label{lem:wheel}
    A $6$-wheel is not separation $2$-choosable.
\end{lemma}

\begin{theorem}
\label{prop:noplanar}
    There is no planar graph $G$ with 
    $\ch_{sep}(G) < \ch_{ad}(G) < \single(G)$.
\end{theorem}

\begin{proof}
    Assume, for a contradiction, that there is such a graph $G$. 
    Then $\ch_{sep}(G)=2$
    and $\single(G) =4$. We assume that $G$ is vertex-minimal with 
    respect to the property that $|\single(G) - \ch_{sep}(G)| \geq 2$,
    so $G$ is connected. Moreover, for every vertex $v$ of $G$, 
    $\single(G-v)=3$, and so, $\delta(G) \geq 3$.
	%This implies that $\delta(G)=3$.
	%and so $G$ is not outerplanar, 
	%and thus contains a subdivision of $K_4$ or of $K_{2,3}$.
Additionally, 
    by Proposition \ref{prop:maxdeg}, $\Delta(G) \geq 5$.

By Proposition \ref{prop:disjointcycles},
$G$ does not contain two $\geq 4$-cycles with 
at most one common vertex. 
Since $G$ has minimum degree at least $3$, it follows that it 
is $2$-connected. (Since otherwise, if $v$ is a cut-vertex, then one of 
the components of $G-v$ is outerplanar, because its longest cycle is a 
triangle.)




Let $x$ be a vertex of degree at least $5$ in $G$, and consider the set 
of neighbors $N_G(x)$ of $x$ in $G$. We shall assume that $x$ has
degree $5$ in $G$; the case when $d_G(x) > 5$ can be dealt with 
similarly.


Now, since $G$ is $2$-connected and $\delta(G-x) \geq 2$,
$G-x$ contains a cycle. Suppose first that none of the vertices
in $N_G(x)$ is contained in a cycle of $G-x$, and let $T$ be a 
vertex-minimal subtree of $G-x$ containing $N_G(x)$ (such a tree exists,
since $G-x$ is connected). Then some vertex $y$ of $N_G(x)$ is a leaf of 
$T$, and since $\delta(G) \geq 3$, this means that $y$ is a cut-vertex 
of $G$, which contradicts that $G$ is $2$-connected. 
%otherwise there would be path from a neighbor of $y$ in 
%$G-V(T)$ to a vertex of $N_G(x)$
Thus, some vertex of $N_G(x)$ is contained in a cycle in $G-x$.

Let $C$ be a cycle in $G-x$ containing the largest number of vertices
from $N_G(x)$. 
Suppose first that all vertices of $N_G(x)$ are contained in $C$.
If $C$ is a $5$-cycle, then by Lemma \ref{lem:wheel}, 
$G$ is not separation
$2$-choosable, a contradiction. On the other hand, if $C$ contains at least $6$ vertices,
then the same conclusion follows from Proposition \ref{prop:disjointcycles}.

Suppose now that four vertices from $N_G(x)$ are contained in $C$,
and that $x_5 \in N_G(x)$ is not contained in $C$.
Since $G-x$ is connected, there is a shortest path $P$ in $G-x$ from
$x_5$ to $V(C)$. If the endpoint of $P$ on $V(C)$ is not in $N_G(x)$,
or $P$ has at least three vertices, then Proposition \ref{prop:disjointcycles}
implies that $G$ is not separation $2$-choosable, so we assume
that $x_5$ is adjacent to $x_4 \in V(C) \cap N_G(x)$. 
By the choice of $C$, 
$x_5$ has no other neighbor on $C$, so since $\delta(G) \geq 3$,
$x_5$ has a neighbor $u \notin N_G(x)$. Again, by the choice
of $C$, if there is a path from $u$ to $V(C)$ in $G-x_5$, then
$x_4$ is the endpoint of such a path. However, then $G$ contains two $\geq 4$-cycles with
at most one common vertex, a contradiction to the separation $2$-choosability of $G$.
On the other hand, if there is no such path from $x_5$ to $V(C)$,
then $x_5$ is a cut-vertex of $G$, a contradiction.

Now assume at most three vertices from $N_G(x)$ are contained in $C$,
and let $x_4,x_5$ be vertices in $N_G(x)$ that are not contained in
$C$. Suppose first that $x_4$ and $C$ are contained in the same
block in $G-x$. Then there must be two internally disjoint paths from $x_4$ to $V(C))$ in $G-x$
with distinct endpoints in $C$. Since $G$ is planar,
this contradicts the choice
of $C$. Thus, $x_4$ and $C$ are contained in different blocks in $G-x$,
as is also $x_5$ and $C$. 

Consider the block graph of $G-x$. %Definition from Diestel
Now, if $x_4$ and $x_5$ are contained in the same block in $G-x$, then
either they are contained in a common cycle or, if $x_4x_5$
is a cut-edge of $G-x$, then since $\delta(G) \geq 3$, 
there is an end-block
in $G$ that contains a $\geq 4$-cycle that is disjoint from $C$. 
In both cases, $G$ contains 
two $\geq 4$-cycles with at most common vertex, 
contradicting that $G$ is
separation $2$-choosable. Hence, $x_4$, $x_5$ and $C$ are contained in
three different blocks in $G - x$. Now, since $\delta(G) \geq 3$, this 
implies that an end-block of $G-x$ contains a 
$\geq 4$-cycle with at most one
common vertex with $C$, a contradiction to the fact that $G$ is
separation $2$-choosable.
\end{proof}

\end{document}
